\Needspace{10\baselineskip}\subsection*{CHEM: inferences}
\addcontentsline{toc}{subsection}{CHEM: inferences}
\begin{inftable}
CH.I01 & Combustion releases phlogiston into the air. & Interpretation: a substance present in every combustible leaves it and passes into the air when it burns. No such substance exists; the air instead loses oxygen. & CH08$^{\circ}$, CH01a$^{\ast}$, CH11$^{\circ}$ & \textbf{U} \\
CH.I02 & Combustion in a closed volume of air stops when the air is saturated with phlogiston. & Interpretation: the air can take up only a limited amount of a substance given off by the burning body. In T2 combustion stops because the oxygen of the enclosed air is used up. & CH08$^{\circ}$, CH01d$^{\ast}$, CH03b$^{\ast}$ & \textbf{U} \\
CH.I03 & Combustion requires air. & Combustion requires air (more exactly, the oxygen in it). & CH08$^{\circ}$, CH11$^{\circ}$ & \textbf{H} \\
CH.I04 & A metal is a compound of its calx and phlogiston. & Interpretation: the metal is its oxide plus a further constituent. T2 has no such constituent; the oxide is the metal plus oxygen. & CH07$^{\circ}$, CH06$^{\circ}$, CH01a$^{\ast}$ & \textbf{U} \\
CH.I05 & Calcination of a metal requires air. & Oxidation of a metal by heating requires air. & CH09, CH07$^{\circ}$, CH11$^{\circ}$ & \textbf{H} \\
CH.I06 & The calx weighs more than the metal from which it was made. & The oxide weighs more than the metal from which it was made. & CH06$^{\circ}$, CH07$^{\circ}$, CH09 & \textbf{H} \\
CH.I07 & Heating a calx with charcoal restores the metal. & Heating an oxide with carbon reduces it to the metal. & CH06$^{\circ}$, CH20$^{\circ}$, CH10$^{\circ}$, CH07$^{\circ}$ & \textbf{H} \\
CH.I08 & In reduction, the calx takes phlogiston from the charcoal. & Interpretation: the oxide gains something from the carbon. In T2 the carbon gains oxygen from the oxide. & CH10$^{\circ}$, CH06$^{\circ}$, CH20$^{\circ}$, CH01b$^{\ast}$ & \textbf{U} \\
CH.I09 & Charcoal burns almost without residue because it is nearly pure phlogiston. & Interpretation: charcoal is almost wholly the principle of inflammability. In T2 charcoal is nearly pure carbon, which burns to a gas. & CH20$^{\circ}$, CH01a$^{\ast}$, CH08$^{\circ}$ & \textbf{U} \\
CH.I10 & The red calx of mercury heated alone yields the metal and dephlogisticated air. & Mercuric oxide heated alone yields mercury and oxygen. & CH06$^{\circ}$, CH07$^{\circ}$, CH02a & \textbf{H} \\
CH.I11 & Dephlogisticated air supports combustion and respiration longer than common air. & Oxygen supports combustion and respiration longer than atmospheric air. & CH02a, CH08$^{\circ}$, CH23, CH11$^{\circ}$ & \textbf{H} \\
CH.I12 & Dephlogisticated air is avid of phlogiston because it has been deprived of it. & Interpretation: the new air is common air emptied of a substance it therefore seeks. In T2 it is oxygen, the active part of air, not a residue. & CH02b$^{\ast}$, CH01a$^{\ast}$ & \textbf{U} \\
CH.I13 & Respiration and combustion both render air unfit for further combustion. & Respiration and combustion both render air unfit for further combustion (both consume its oxygen). & CH23, CH08$^{\circ}$, CH11$^{\circ}$ & \textbf{H} \\
CH.I14 & Air in which a candle has burned out, washed with lime water, extinguishes flame and kills animals. & Nitrogen, the residue after combustion and absorption of carbon dioxide, extinguishes flame and kills animals. & CH03a, CH08$^{\circ}$, CH05 & \textbf{H} \\
CH.I15 & Fixed air is released when chalk is heated or treated with acids. & Carbon dioxide is released when calcium carbonate is heated or treated with acids. & CH05 & \textbf{H} \\
CH.I16 & Fixed air is produced by respiration and by burning charcoal. & Carbon dioxide is produced by respiration and by burning carbon. & CH05, CH23, CH20$^{\circ}$, CH08$^{\circ}$ & \textbf{H} \\
CH.I17 & Inflammable air is obtained by dissolving iron or zinc in dilute vitriolic acid. & Hydrogen is obtained by dissolving iron or zinc in dilute sulphuric acid. & CH04a, CH07$^{\circ}$, CH18$^{\circ}$ & \textbf{H} \\
CH.I18 & Inflammable air exploded with dephlogisticated air leaves water, and nearly the whole of both airs is consumed. & Hydrogen exploded with oxygen yields water, nearly the whole of both gases being consumed. & CH04a, CH02c, CH21$^{\circ}$ & \textbf{H} \\
CH.I19 & Water is an element. & Water is a simple substance. [Denied: water is a compound of hydrogen and oxygen.] & CH21$^{\circ}$, CH22$^{\circ}$ & \textbf{X} \\
CH.I20 & Sulphur is composed of vitriolic acid and phlogiston. & Interpretation: sulphur is sulphuric acid combined with the principle of inflammability. In T2 sulphur is simple and sulphuric acid contains sulphur and oxygen. & CH19$^{\circ}$, CH18$^{\circ}$, CH01a$^{\ast}$ & \textbf{U} \\
CH.I21 & The goodness of an air is measured by its diminution with nitrous air over water. & The oxygen content of an air is measured by its contraction with nitric oxide over water. & CH12$^{\circ}$, CH13 & \textbf{H} \\
CH.I22 & The goodness of an air is its capacity to receive phlogiston. & Interpretation: an air is good in proportion as it can absorb what burning bodies give off. In T2 it is good in proportion to its oxygen. & CH12$^{\circ}$, CH01d$^{\ast}$ & \textbf{U} \\
CH.I23 & Dephlogisticated nitrous air supports the flame of a candle. & Nitrous oxide supports the flame of a candle. & CH14, CH08$^{\circ}$ & \textbf{H} \\
CH.I24 & Marine acid air and alkaline air combine to form a white salt (sal ammoniac). & Hydrogen chloride and ammonia combine to form ammonium chloride. & CH15, CH17 & \textbf{H} \\
CH.I25 & Dephlogisticated marine acid bleaches vegetable colours. & Chlorine (oxygenated muriatic acid) bleaches vegetable colours. & CH16 & \textbf{H} \\
CH.I26 & Metals resemble one another because they all contain phlogiston. & Interpretation: the metals share a common constituent that explains their resemblance. T2 recognizes no such constituent and in 1789 had no explanation of the resemblance. & CH07$^{\circ}$, CH01a$^{\ast}$ & \textbf{U} \\
CH.I27 & Air is an element. & Atmospheric air is a simple substance. [Denied: it is a mixture of oxygen and nitrogen.] & CH11$^{\circ}$, CH22$^{\circ}$ & \textbf{X} \\
CH.I28 & Inflammable air is phlogiston, or nearly so. & Hydrogen is hydrogen, or nearly so (true on this reading: the token of phlogiston refers to hydrogen). & CH04a, CH01c & \textbf{H} \\
CH.I29 & Marsh air burns with a flame. & Methane burns with a flame. & CH04b & \textbf{H} \\
CH.I30 & Heavy inflammable air is produced when charcoal is heated with a metallic calx. & Carbon monoxide is produced when carbon is heated with a metal oxide. & CH04c, CH20$^{\circ}$, CH06$^{\circ}$ & \textbf{H} \\
CH.I31 & Sulphur burned in air yields an acid with water. & Sulphur burned in air yields an acid with water. & CH19$^{\circ}$, CH08$^{\circ}$, CH21$^{\circ}$ & \textbf{H} \\
\bottomrule
\end{inftable}

\Needspace{10\baselineskip}\subsection*{ASTR: inferences}
\addcontentsline{toc}{subsection}{ASTR: inferences}
\begin{inftable}
AS.I01 & Each planet moves on an epicycle whose centre moves on a deferent. & The geocentric motion of a planet is the composition of its heliocentric orbit and the Earth's orbit. [Holds to first order: circular approximations to ellipses.] & AS01a, AS06a, AS05a & \textbf{R} \\
AS.I02 & A superior planet retrogrades when it is in opposition to the Sun. & A superior planet appears to retrograde when it is in opposition to the Sun. & AS11$^{\circ}$, AS01a, AS13$^{\circ}$ & \textbf{H} \\
AS.I03 & A superior planet is brightest when it retrogrades. & A superior planet is brightest when it appears to retrograde (it is then nearest the Earth). & AS11$^{\circ}$, AS01a & \textbf{H} \\
AS.I04 & The epicycle centre of Venus and Mercury always lies on the line from the Earth to the mean Sun. & The Sun, centre of the orbits of Venus and Mercury, lies on the line from the Earth to the Sun. [Holds exactly with the true Sun; the mean Sun is a first-order substitute.] & AS06b, AS05b, AS01d, AS13$^{\circ}$, AS12$^{\circ}$ & \textbf{R} \\
AS.I05 & The radius of a superior planet's epicycle stays parallel to the line from the Earth to the mean Sun. & The component of a superior planet's geocentric motion due to the Earth's orbit stays parallel to the Earth--Sun line. [Holds with circular orbits.] & AS06a, AS01a, AS13$^{\circ}$, AS12$^{\circ}$ & \textbf{R} \\
AS.I06 & The Earth is at rest. & The Earth is at rest. [Denied: it rotates daily and revolves annually.] & AS12$^{\circ}$ & \textbf{X} \\
AS.I07 & Heavy bodies move naturally toward the centre of the universe. & Heavy bodies move naturally toward the centre of the Earth. & AS15a & \textbf{H} \\
AS.I08 & The primum mobile communicates the daily motion to the inner orbs. & Interpretation: an outermost body turns daily and drags the inner heavens with it. T2 recognizes no such body; the daily motion is the Earth's rotation. & AS03a$^{\ast}$, AS16$^{\circ}$, AS04a$^{\ast}$ & \textbf{U} \\
AS.I09 & All celestial bodies share a westward diurnal motion with a period of one sidereal day. & All celestial bodies appear to move westward once each sidereal day. & AS16$^{\circ}$, AS02$^{\circ}$, AS01a & \textbf{H} \\
AS.I10 & The diurnal motion is a real rotation of the heavens about the centre of the universe. & Interpretation: the heavens physically rotate about a fixed central point. In T2 nothing rotates about such a point; the Earth rotates. & AS16$^{\circ}$, AS15b$^{\ast}$ & \textbf{U} \\
AS.I11 & The planets are carried by solid orbs that nest without gaps. & Interpretation: solid shells of a celestial matter carry the planets. T2 has no such shells; Kepler sought instead a force from the Sun. & AS04a$^{\ast}$, AS01a & \textbf{U} \\
AS.I12 & The heavens consist of aether, which is ungenerated and unalterable. & Interpretation: the heavens are made of a special unalterable matter. T2 (with Tycho's parallaxes of the nova of 1572 and the comet of 1577) recognizes no such matter. & AS24a$^{\ast}$ & \textbf{U} \\
AS.I13 & The fixed stars keep their mutual positions throughout the year. & The stars keep their mutual positions throughout the year (no annual parallax is detectable). & AS02$^{\circ}$ & \textbf{H} \\
AS.I14 & Precession is a slow eastward motion of the sphere of the fixed stars. & Precession is a slow motion of the stars relative to the equinoxes. [Denied: it is a slow motion of the Earth's axis.] & AS19$^{\circ}$, AS02$^{\circ}$ & \textbf{X} \\
AS.I15 & The epicycle centre moves uniformly as seen from the equant. & The planet moves uniformly in angle about the empty focus of its ellipse. [Holds to first order in the eccentricity.] & AS08a, AS06a & \textbf{R} \\
AS.I16 & The equant lies at twice the distance of the deferent's centre from the Earth (bisected eccentricity). & The empty focus lies at twice the distance of the ellipse's centre from the Sun. [Holds: Kepler confirmed the bisection of the eccentricity.] & AS08a, AS07$^{\circ}$ & \textbf{R} \\
AS.I17 & The period of the second anomaly is the synodic period. & The period of the apparent effect of the Earth's motion is the synodic period. & AS10, AS20 & \textbf{H} \\
AS.I18 & The first anomaly depends on the planet's longitude in the zodiac. & The non-uniformity of the planet's orbital motion depends on its heliocentric longitude. & AS09 & \textbf{H} \\
AS.I19 & The mean longitude of Venus and Mercury advances with a period of one year. & The mean geocentric longitude of Venus and Mercury advances with the Earth's sidereal period of one year. & AS21b, AS01d & \textbf{H} \\
AS.I20 & The order of the planets from the Earth is Moon, Mercury, Venus, Sun, Mars, Jupiter, Saturn. & The order of distances from the Earth is Moon, Mercury, Venus, Sun, Mars, Jupiter, Saturn. [Denied: Venus and Mercury are sometimes farther than the Sun; the order about the Sun is Mercury, Venus, Earth, Mars, Jupiter, Saturn.] & AS22$^{\circ}$, AS01b, AS14$^{\circ}$, AS13$^{\circ}$ & \textbf{X} \\
AS.I21 & Saturn completes its circuit of the zodiac in about 29 1/2 years. & Saturn completes its orbit in about 29 1/2 years. & AS21a, AS01a & \textbf{H} \\
AS.I22 & Venus never departs from the Sun by more than about 47 degrees. & Venus never appears more than about 47 degrees from the Sun. & AS23, AS01d, AS13$^{\circ}$ & \textbf{H} \\
AS.I23 & Comets are exhalations below the Moon. & Comets are below the Moon. [Denied on Tycho's parallax of the comet of 1577.] & AS25$^{\circ}$, AS14$^{\circ}$ & \textbf{X} \\
AS.I24 & The tropical year is shorter than the time of the Sun's return to the same fixed star. & The tropical year is shorter than the sidereal year, the difference being due to precession. & AS18, AS19$^{\circ}$, AS13$^{\circ}$ & \textbf{H} \\
AS.I25 & The Sun moves uniformly on an eccentric circle about the Earth. & The Earth moves uniformly on an eccentric circle about the Sun. [Holds to first order; Kepler: an ellipse with the area law.] & AS13$^{\circ}$, AS07$^{\circ}$, AS12$^{\circ}$ & \textbf{R} \\
AS.I26 & The Moon is the planet nearest the Earth. & The Moon is the celestial body nearest the Earth. & AS14$^{\circ}$, AS01c & \textbf{H} \\
AS.I27 & The Sun moves annually along the ecliptic. & The Sun appears to move annually along the ecliptic. & AS13$^{\circ}$, AS17 & \textbf{H} \\
AS.I28 & The Moon's motion has an inequality depending on its position in its own circuit. & The Moon's motion has an inequality depending on its position in its orbit (the equation of centre). & AS06c, AS14$^{\circ}$ & \textbf{H} \\
AS.I29 & The orbs fix the absolute distances of the planets. & Interpretation: nested solid shells fix how far each planet is. T2 fixes relative distances by the geometry of the orbits instead. & AS04b$^{\ast}$ & \textbf{U} \\
AS.I30 & Retrograde motion is a real motion of the planet. & Retrograde motion is a real motion of the planet. [Denied: it is apparent.] & AS11$^{\circ}$ & \textbf{X} \\
AS.I31 & The second anomaly of a superior planet is completed when it returns to the same elongation from the Sun. & The synodic cycle of a superior planet is completed when it returns to the same elongation from the Sun. & AS10, AS23 & \textbf{H} \\
AS.I32 & Eclipses occur only when the Moon is near the ecliptic. & Eclipses occur only when the Moon is near the ecliptic. & AS17, AS14$^{\circ}$ & \textbf{H} \\
AS.I33 & The orb of Mars turns physically uniformly about the equant. & Interpretation: a solid shell rotates uniformly about an off-centre point. T2 has no shells; Kepler took the non-uniformity to be real. & AS08b$^{\ast}$, AS04a$^{\ast}$ & \textbf{U} \\
AS.I34 & The Sun is a planet like the others. & The Sun is a planet like the others. [Denied: it is the centre about which the planets move.] & AS13$^{\circ}$, AS01b & \textbf{X} \\
\bottomrule
\end{inftable}

\Needspace{10\baselineskip}\subsection*{MECH: inferences}
\addcontentsline{toc}{subsection}{MECH: inferences}
\begin{inftable}
ME.I01 & The mass of a body is the same in every frame. & The mass of a body is the same in every frame. [Holds on the rest-mass reading; on the relativistic-mass reading only for $v \ll c$.] & ME01$^{\circ}$ & \textbf{R} \\
ME.I02 & The mass of a composite is the sum of the masses of its parts. & The mass of a composite is the sum of the masses of its parts. [Holds on the relativistic-mass reading; for rest mass only when binding and kinetic energies are negligible.] & ME01$^{\circ}$ & \textbf{R} \\
ME.I03 & Mass is conserved in all interactions. & Mass is conserved in all interactions. [Holds on the relativistic-mass reading; rest mass is not conserved.] & ME01$^{\circ}$ & \textbf{R} \\
ME.I04 & Momentum is the product of mass and velocity. & Momentum is the product of mass and velocity. [Holds exactly on the relativistic-mass reading; on the rest-mass reading only for $v \ll c$.] & ME03$^{\circ}$, ME01$^{\circ}$, ME05$^{\circ}$ & \textbf{R} \\
ME.I05 & Momentum is conserved in an isolated system. & Relativistic momentum is conserved in an isolated system. & ME03$^{\circ}$ & \textbf{H} \\
ME.I06 & Force equals the rate of change of momentum. & The force on a charge equals the rate of change of its relativistic momentum. & ME02a, ME03$^{\circ}$ & \textbf{H} \\
ME.I07 & The gravitational force between two bodies varies as the product of their masses and inversely as the square of their distance. & Interpretation: a force reaches from the Sun to each planet in proportion to their masses and inversely as the square of the distance. In general relativity there is no such force; the Newtonian law is recovered as the weak-field limit of geodesic motion (ME.I08). & ME02b$^{\ast}$, ME01$^{\circ}$, ME11$^{\circ}$ & \textbf{U} \\
ME.I08 & In a gravitational field a body accelerates down the gradient of the potential. & In a weak static field a slow body accelerates down the gradient of $(g_{00}+1)c^2/2$. [Weak-field, slow-motion limit.] & ME12$^{\circ}$, ME21$^{\circ}$ & \textbf{R} \\
ME.I09 & Simultaneity of distant events is the same for all observers. & Simultaneity of distant events is the same in all inertial frames. [Holds as $v\,\Delta x/c^2 \to 0$.] & ME08$^{\circ}$ & \textbf{R} \\
ME.I10 & Velocities compose by addition. & Velocities compose by addition. [Holds for $v \ll c$; in general by the relativistic formula.] & ME05$^{\circ}$, ME10$^{\circ}$ & \textbf{R} \\
ME.I11 & Inertial frames are related by Galilean transformations. & Inertial frames are related by Galilean transformations. [Limit of the Lorentz transformation as $v/c \to 0$.] & ME09$^{\circ}$, ME10$^{\circ}$ & \textbf{R} \\
ME.I12 & The laws of mechanics are the same in all inertial frames. & The laws of physics are the same in all inertial frames. & ME09$^{\circ}$ & \textbf{H} \\
ME.I13 & Absolute space remains immovable and is not acted on by bodies. & Interpretation: a background space exists that is immovable and is not acted on by bodies. T2 recognizes no state of absolute rest; in general relativity the geometry is itself dynamical. & ME06a$^{\ast}$ & \textbf{U} \\
ME.I14 & The concavity of the water in the rotating bucket reveals its true rotation. & The concavity of the water reveals rotation relative to the local inertial structure. & ME06b & \textbf{H} \\
ME.I15 & Absolute time flows equably for all observers. & Interpretation: one time flows at the same rate for every observer. T2 has proper times along worldlines and frame-relative coordinate times. & ME07a$^{\ast}$ & \textbf{U} \\
ME.I16 & Clocks once synchronized remain synchronized whatever their relative motion. & Clocks once synchronized remain synchronized whatever their relative motion. [Holds as $v/c \to 0$ and in weak fields.] & ME07b, ME08$^{\circ}$ & \textbf{R} \\
ME.I17 & The length of a body does not depend on its motion. & The length of a body does not depend on its motion. [Holds as $v/c \to 0$.] & ME14$^{\circ}$ & \textbf{R} \\
ME.I18 & Light travels with a fixed speed relative to the ether. & Interpretation: there is a medium relative to which light has its speed. T2 has no such medium; light has speed $c$ in every inertial frame. & ME15$^{\circ}$, ME16a$^{\ast}$ & \textbf{U} \\
ME.I19 & The measured speed of light depends on the observer's motion. & The measured speed of light depends on the observer's motion. [Denied: it is $c$ in every inertial frame.] & ME15$^{\circ}$, ME05$^{\circ}$ & \textbf{X} \\
ME.I20 & Electromagnetic disturbances in the ether propagate with the speed of light. & Electromagnetic disturbances in the field propagate with speed $c$. & ME16b, ME15$^{\circ}$ & \textbf{H} \\
ME.I21 & A planet attracted by the Sun alone describes a fixed ellipse. & A planet orbiting the Sun alone follows a fixed ellipse. [Holds in the weak-field limit; general relativity adds a perihelion advance.] & ME13$^{\circ}$, ME11$^{\circ}$ & \textbf{R} \\
ME.I22 & Gravitational and inertial mass are proportional. & Gravitational and inertial mass are proportional (the equivalence principle). & ME01$^{\circ}$, ME11$^{\circ}$ & \textbf{H} \\
ME.I23 & Gravitation acts instantaneously across any distance. & Interpretation: gravitation acts at once across any distance. In T2 changes in the gravitational field propagate at the speed of light. & ME23a$^{\ast}$, ME11$^{\circ}$ & \textbf{U} \\
ME.I24 & An impulse given to one end of a rigid body moves the other end at once. & Interpretation: a body could transmit a push instantaneously. T2 excludes signals faster than light, and so excludes perfectly rigid bodies. & ME22a$^{\ast}$ & \textbf{U} \\
ME.I25 & A rod moved slowly from one place to another keeps its length. & A rod moved slowly from place to place keeps its proper length. & ME22b, ME14$^{\circ}$ & \textbf{H} \\
ME.I26 & Kinetic energy is one half of mass times velocity squared. & Kinetic energy is one half of mass times velocity squared. [Leading term of $(\gamma-1)m_0c^2$ for $v \ll c$.] & ME04$^{\circ}$, ME01$^{\circ}$, ME05$^{\circ}$ & \textbf{R} \\
ME.I27 & Energy is conserved in an isolated system. & Energy is conserved in an isolated system. & ME20 & \textbf{H} \\
ME.I28 & The weight of a body is the gravitational force on it. & Interpretation: weight is the pull of gravity on a body. In general relativity no force acts on a freely falling body; what a balance registers is the support force. & ME19a$^{\ast}$, ME02b$^{\ast}$ & \textbf{U} \\
ME.I29 & The centre of mass of an isolated system moves uniformly. & The centre of energy of an isolated system moves uniformly in any inertial frame. & ME18 & \textbf{H} \\
ME.I30 & Acceleration is the same in all inertial frames. & Acceleration is the same in all inertial frames. [Holds as $v/c \to 0$.] & ME21$^{\circ}$, ME09$^{\circ}$ & \textbf{R} \\
ME.I31 & The strength of gravitation is fixed by one universal constant. & The strength of gravitation is fixed by one universal constant $G$. & ME17, ME11$^{\circ}$ & \textbf{H} \\
ME.I32 & A spring balance holding a body at rest on the Earth registers its weight. & A spring balance holding a body at rest on the Earth registers the support force that keeps it from free fall. & ME19b & \textbf{H} \\
ME.I33 & The Earth's motion through the ether should alter the round-trip time of light along perpendicular arms. & Interpretation: the Earth moves through a medium, which should alter the round-trip time of light. In T2 there is no medium to move through, and no alteration is expected. & ME16c$^{\ast}$, ME15$^{\circ}$ & \textbf{U} \\
ME.I34 & Bodies falling freely together show no relative acceleration in a uniform field. & Bodies falling freely together show no relative acceleration in a uniform field. & ME09$^{\circ}$, ME21$^{\circ}$ & \textbf{H} \\
ME.I35 & The centre of the world is at rest in absolute space. & Interpretation: the centre of the solar system is at absolute rest. T2 has no absolute rest. & ME06c$^{\ast}$ & \textbf{U} \\
\bottomrule
\end{inftable}

\Needspace{10\baselineskip}\subsection*{HEAT: inferences}
\addcontentsline{toc}{subsection}{HEAT: inferences}
\begin{inftable}
HE.I01 & Caloric is a subtle fluid whose particles repel one another and are attracted by ordinary matter. & Interpretation: heat is a self-repelling fluid attracted by matter. T2 has no such fluid; heat is energy transferred. & HE01a$^{\ast}$ & \textbf{U} \\
HE.I02 & The total quantity of caloric is conserved. & Interpretation: the quantity of a heat-substance is fixed. In T2 energy is conserved but heat is not: it is converted to and from work. & HE01c$^{\ast}$ & \textbf{U} \\
HE.I03 & When hot and cold water are mixed, the heat lost by the one equals the heat gained by the other. & When hot and cold water are mixed (no work done), the heat lost by the one equals the heat gained by the other. & HE02$^{\circ}$, HE03$^{\circ}$ & \textbf{H} \\
HE.I04 & Heat passes spontaneously from hotter to colder bodies. & Heat passes spontaneously from hotter to colder bodies. & HE02$^{\circ}$, HE03$^{\circ}$ & \textbf{H} \\
HE.I05 & Ice at its melting point absorbs heat without rising in temperature. & Ice at its melting point absorbs heat without rising in temperature. & HE04$^{\circ}$, HE03$^{\circ}$ & \textbf{H} \\
HE.I06 & Latent heat is caloric combined with the substance. & Interpretation: the absorbed heat is held combined in the substance. In T2 it goes into the molecular potential energy of the new phase. & HE04$^{\circ}$, HE07a$^{\ast}$ & \textbf{U} \\
HE.I07 & The temperature rise of a body is the heat added divided by its heat capacity. & The temperature rise of a body is the heat added divided by its heat capacity. & HE05$^{\circ}$, HE02$^{\circ}$, HE03$^{\circ}$ & \textbf{H} \\
HE.I08 & Specific heat is the capacity of a body for caloric. & Interpretation: specific heat measures how much heat-fluid a body can hold. In T2 it is the energy needed per degree per unit mass. & HE05$^{\circ}$, HE01a$^{\ast}$ & \textbf{U} \\
HE.I09 & Air heats on sudden compression and cools on sudden expansion. & Air heats on sudden compression and cools on sudden expansion. & HE10, HE18$^{\circ}$, HE03$^{\circ}$ & \textbf{H} \\
HE.I10 & Sound compresses air adiabatically, so its speed is $\sqrt{\gamma p/\rho}$. & Sound compresses air adiabatically, so its speed is $\sqrt{\gamma p/\rho}$. & HE12, HE11, HE10 & \textbf{H} \\
HE.I11 & The specific heat of a gas at constant pressure exceeds that at constant volume. & The specific heat of a gas at constant pressure exceeds that at constant volume. & HE05$^{\circ}$, HE11, HE18$^{\circ}$ & \textbf{H} \\
HE.I12 & The efficiency of a reversible engine depends only on the temperatures of its reservoirs. & The efficiency of a reversible engine depends only on the temperatures of its reservoirs. & HE13$^{\circ}$, HE14, HE03$^{\circ}$ & \textbf{H} \\
HE.I13 & In a cycle, the heat given to the cold reservoir equals the heat taken from the hot. & In a cycle, the heat given to the cold reservoir equals the heat taken from the hot. [Denied: it is less by the work done.] & HE02$^{\circ}$, HE13$^{\circ}$ & \textbf{X} \\
HE.I14 & Motive power is produced by the fall of caloric through a temperature difference, as water falls through a wheel. & Work is produced by the passage of entropy through a temperature difference, $W = S(T_1 - T_2)$ for a reversible engine. [Holds on Callendar's reading of the token; fails on Kelvin's.] & HE01d, HE13$^{\circ}$, HE03$^{\circ}$ & \textbf{R} \\
HE.I15 & Frictional heat is caloric released from bodies whose capacity is reduced by abrasion. & Interpretation: frictional heat is heat-fluid squeezed out of abraded bodies. In T2 it is work converted into heat. & HE15$^{\circ}$, HE05$^{\circ}$, HE01a$^{\ast}$ & \textbf{U} \\
HE.I16 & Heat is conducted with a flux proportional to the temperature gradient. & Heat is conducted with a flux proportional to the temperature gradient. & HE08, HE03$^{\circ}$ & \textbf{H} \\
HE.I17 & All bodies emit radiant heat at all temperatures, and an enclosure comes to equilibrium by exchange. & All bodies emit thermal radiation at all temperatures, and an enclosure comes to equilibrium by exchange. & HE09$^{\circ}$, HE16 & \textbf{H} \\
HE.I18 & Radiant heat is caloric in free flight. & Interpretation: radiant heat is the heat-fluid travelling freely. In T2 it is electromagnetic radiation. & HE09$^{\circ}$, HE01a$^{\ast}$ & \textbf{U} \\
HE.I19 & Bodies in contact come to a common temperature. & Bodies in contact come to a common temperature. & HE16, HE03$^{\circ}$ & \textbf{H} \\
HE.I20 & Gases at constant pressure expand by the same fraction for the same rise in temperature. & Gases at constant pressure expand by the same fraction for the same rise in temperature. [Holds for ideal gases.] & HE17$^{\circ}$, HE18$^{\circ}$, HE03$^{\circ}$ & \textbf{R} \\
HE.I21 & The elasticity of a gas arises from static repulsion between caloric atmospheres. & Interpretation: the spring of a gas comes from static repulsion of heat-fluid atmospheres. In T2 pressure comes from molecular impacts. & HE19$^{\circ}$, HE01a$^{\ast}$, HE18$^{\circ}$ & \textbf{U} \\
HE.I22 & At constant temperature, the elasticity of a gas varies inversely as its volume. & At constant temperature the pressure of a gas varies inversely as its volume. [Holds for ideal gases.] & HE19$^{\circ}$, HE18$^{\circ}$ & \textbf{R} \\
HE.I23 & Respiration is a slow combustion and the source of animal heat. & Respiration is a slow combustion and the source of body heat. & HE23, HE02$^{\circ}$ & \textbf{H} \\
HE.I24 & The ice calorimeter measures heat by the quantity of ice melted. & The ice calorimeter measures heat by the quantity of ice melted. & HE21, HE02$^{\circ}$, HE04$^{\circ}$ & \textbf{H} \\
HE.I25 & Evaporation absorbs heat. & Evaporation absorbs heat. & HE20$^{\circ}$, HE04$^{\circ}$ & \textbf{H} \\
HE.I26 & The thermometer measures the quantity of free caloric in a body. & Interpretation: the thermometer reads the density of uncombined heat-fluid. In T2 it reads temperature, which is not a quantity of anything. & HE22$^{\circ}$, HE06a$^{\ast}$ & \textbf{U} \\
HE.I27 & In Laplace's calculation the caloric released by compression stays in the compressed parcel. & In Laplace's calculation the heat released by compression stays in the compressed parcel (the compression is adiabatic). & HE01b, HE10 & \textbf{H} \\
HE.I28 & Heat cannot be generated without limit from a finite body. & Heat cannot be generated without limit from a finite body. [Denied: work can be converted into heat without limit, as in Rumford's boring.] & HE02$^{\circ}$, HE15$^{\circ}$ & \textbf{X} \\
HE.I29 & A vapour is a liquid combined with caloric. & Interpretation: a vapour is a liquid chemically combined with heat-fluid. In T2 it is the gaseous phase. & HE20$^{\circ}$, HE07a$^{\ast}$ & \textbf{U} \\
HE.I30 & Thermometers agree on the order of bodies by hotness. & Thermometers agree on the order of bodies by temperature. & HE22$^{\circ}$, HE03$^{\circ}$ & \textbf{H} \\
\bottomrule
\end{inftable}

\Needspace{10\baselineskip}\subsection*{QUAN: inferences}
\addcontentsline{toc}{subsection}{QUAN: inferences}
\begin{inftable}
QU.I01 & An electron in an atom moves on a definite orbit about the nucleus. & Interpretation: the bound electron follows a closed path about the nucleus. In T2 a stationary state assigns no path to the electron. & QU02a$^{\ast}$, QU06$^{\circ}$, QU07$^{\circ}$, QU08$^{\circ}$ & \textbf{U} \\
QU.I02 & An electron revolving in an atom radiates continuously and spirals into the nucleus. & Interpretation: the bound electron loses energy continuously and falls into the nucleus. In T2 the ground state is stable and does not radiate. & QU02a$^{\ast}$, QU10$^{\circ}$, QU08$^{\circ}$ & \textbf{U} \\
QU.I03 & An accelerated charge radiates with power given by the Larmor formula. & An accelerated charge radiates with power given by the Larmor formula. [Holds for macroscopic currents and free charges; not for stationary states.] & QU10$^{\circ}$, QU19$^{\circ}$ & \textbf{R} \\
QU.I04 & The frequencies emitted by an atom are those of its electrons' orbital motions and their harmonics. & Interpretation: the emitted frequencies are those of periodic electron motions. In T2 they are differences of stationary-state energies divided by $h$. & QU11$^{\circ}$, QU02a$^{\ast}$ & \textbf{U} \\
QU.I05 & A very highly excited electron radiates at its orbital frequency. & A circulating Rydberg wave packet radiates at its orbital frequency. [Holds in the limit of high quantum numbers: Bohr's correspondence principle.] & QU02b, QU11$^{\circ}$ & \textbf{R} \\
QU.I06 & Every particle has at every time a definite position and momentum. & Every particle has at every time a definite value of the position and of the momentum observable. [Denied in the standard formulation.] & QU03$^{\circ}$, QU04$^{\circ}$ & \textbf{X} \\
QU.I07 & Position and momentum can be measured together with arbitrary precision. & Position and momentum can be measured together with arbitrary precision. [Denied: $\Delta x\,\Delta p \geq \hbar/2$.] & QU03$^{\circ}$, QU04$^{\circ}$ & \textbf{X} \\
QU.I08 & The trajectory of a body is fixed by its position and momentum at one time together with the forces. & The centre of a localized state moves as the classical equations say. [Ehrenfest's theorem: holds for well-localized states in slowly varying potentials.] & QU01c, QU03$^{\circ}$, QU04$^{\circ}$ & \textbf{R} \\
QU.I09 & A charged particle in supersaturated vapour leaves a track along its trajectory. & A charged particle in supersaturated vapour leaves a nearly straight track of condensed droplets. & QU01b, QU06$^{\circ}$ & \textbf{H} \\
QU.I10 & The electron in an atom can have any energy in a continuous range. & The energy of a bound electron can take any value in a continuous range. [Denied: bound-state energies are discrete.] & QU05$^{\circ}$, QU06$^{\circ}$, QU07$^{\circ}$ & \textbf{X} \\
QU.I11 & In equilibrium every degree of freedom has mean energy $kT/2$. & In equilibrium each quadratic degree of freedom has mean energy $kT/2$. [Holds for $kT$ large compared with the level spacing.] & QU12$^{\circ}$, QU05$^{\circ}$ & \textbf{R} \\
QU.I12 & The spectral density of cavity radiation is proportional to $\nu^2 T$ (Rayleigh--Jeans). & The spectral density of cavity radiation is proportional to $\nu^2 T$. [Holds for $h\nu \ll kT$.] & QU13$^{\circ}$, QU12$^{\circ}$ & \textbf{R} \\
QU.I13 & The atomic heat of every solid element is about 6 calories per degree. & The molar heat capacity of every solid element is about $3R$. [Holds at temperatures high compared with its characteristic temperature.] & QU14$^{\circ}$, QU12$^{\circ}$ & \textbf{R} \\
QU.I14 & The energy of photoelectrons increases with the intensity of the light. & The energy of photoelectrons increases with the intensity of the light. [Denied: it depends on the frequency; the intensity fixes their number.] & QU15$^{\circ}$, QU09$^{\circ}$ & \textbf{X} \\
QU.I15 & Light shows interference and diffraction. & Electromagnetic radiation shows interference and diffraction. & QU09$^{\circ}$ & \textbf{H} \\
QU.I16 & The energy of light is spread continuously over the wave front. & The energy of radiation is spread continuously over the wave front. [Denied: it is exchanged with matter in quanta $h\nu$.] & QU09$^{\circ}$, QU05$^{\circ}$ & \textbf{X} \\
QU.I17 & Electrons have a definite charge and mass. & Electrons have a definite charge and mass. & QU06$^{\circ}$, QU19$^{\circ}$ & \textbf{H} \\
QU.I18 & Charge is conserved and comes in multiples of $e$. & Charge is conserved and comes in multiples of $e$ (for free particles). & QU19$^{\circ}$ & \textbf{H} \\
QU.I19 & Nearly all the mass and the positive charge of an atom are in a small nucleus. & Nearly all the mass and the positive charge of an atom are in a small nucleus. & QU08$^{\circ}$, QU07$^{\circ}$ & \textbf{H} \\
QU.I20 & Alpha particles scattered by a nucleus follow Rutherford's angular distribution. & Alpha particles scattered by a nucleus follow Rutherford's distribution (which the quantum treatment of Coulomb scattering reproduces exactly). & QU21, QU08$^{\circ}$ & \textbf{H} \\
QU.I21 & X-rays are diffracted by crystals as short electromagnetic waves. & X-rays are diffracted by crystals as short electromagnetic waves. & QU20, QU09$^{\circ}$ & \textbf{H} \\
QU.I22 & The angular momentum of a body can take any value. & The angular momentum of a body can take any value. [Denied: its components have discrete eigenvalues.] & QU16$^{\circ}$ & \textbf{X} \\
QU.I23 & The angular momentum of an isolated system is conserved. & The angular momentum of an isolated system is conserved. & QU16$^{\circ}$ & \textbf{H} \\
QU.I24 & An atomic magnetic moment can point in any direction relative to a field. & An atomic magnetic moment can point in any direction relative to a field. [Denied: Stern and Gerlach found discrete components.] & QU17$^{\circ}$, QU16$^{\circ}$ & \textbf{X} \\
QU.I25 & The state of a system is a point in phase space, evolving deterministically. & Interpretation: the full state is a set of exact positions and momenta, evolving deterministically. In T2 the state is a vector in Hilbert space. & QU18a$^{\ast}$ & \textbf{U} \\
QU.I26 & The macroscopic state of a gas is fixed by its pressure, volume and temperature. & The thermodynamic macrostate of a gas is fixed by its pressure, volume and temperature. & QU18b & \textbf{H} \\
QU.I27 & An antenna current oscillating at a given frequency radiates at that frequency. & An antenna current oscillating at a given frequency radiates at that frequency. & QU10$^{\circ}$, QU09$^{\circ}$ & \textbf{H} \\
QU.I28 & The spectral lines of hydrogen satisfy the Balmer--Rydberg formula. & The transition frequencies of hydrogen satisfy the Balmer--Rydberg formula. & QU11$^{\circ}$, QU07$^{\circ}$ & \textbf{H} \\
QU.I29 & A planet's orbit is a definite path in space. & A planet's orbit is a definite path in space. [Holds for macroscopic bodies to any accuracy that can be tested.] & QU01c & \textbf{R} \\
QU.I30 & The electron in an atom traces a continuous trajectory fixed by the Coulomb force of the nucleus. & Interpretation: the bound electron has a continuous path determined by the Coulomb force. In T2 the Coulomb potential fixes the stationary states, which assign no path. & QU01a$^{\ast}$, QU08$^{\circ}$ & \textbf{U} \\
\bottomrule
\end{inftable}

\Needspace{10\baselineskip}\subsection*{BIOL: inferences}
\addcontentsline{toc}{subsection}{BIOL: inferences}
\begin{inftable}
BI.I01 & Species are immutable. & Species do not change. [Denied: species descend with modification.] & BI01$^{\circ}$ & \textbf{X} \\
BI.I02 & Varieties differ from species in not being permanent. & Varieties differ from species in not being permanent. [Denied: varieties are incipient species, and the difference is one of degree.] & BI02$^{\circ}$, BI01$^{\circ}$ & \textbf{X} \\
BI.I03 & The variation of a species is confined within fixed limits about its type. & The heritable variation within a species is confined within fixed limits. [Denied.] & BI23$^{\circ}$, BI01$^{\circ}$ & \textbf{X} \\
BI.I04 & The natural system arranges organisms in groups subordinate to groups. & Genealogical classification arranges organisms in groups subordinate to groups. & BI04$^{\circ}$, BI03 & \textbf{H} \\
BI.I05 & The natural system expresses the plan of the Creator. & Interpretation: the hierarchy of groups expresses a design. In T2 it expresses community of descent. & BI04$^{\circ}$, BI05a$^{\ast}$ & \textbf{U} \\
BI.I06 & The members of each great class share a unity of plan. & The members of each great class share an ancestral pattern of structure. & BI05b, BI04$^{\circ}$ & \textbf{H} \\
BI.I07 & The archetype is an idea in the divine mind. & Interpretation: the pattern is an idea in the mind of a designer. In T2 it is the structure of an ancestor. & BI06a$^{\ast}$ & \textbf{U} \\
BI.I08 & All vertebrate skeletons are modifications of one archetypal pattern. & All vertebrate skeletons are modifications of the skeleton of their common ancestor. & BI06b, BI07 & \textbf{H} \\
BI.I09 & The forelimbs of man, bat and porpoise are homologous. & The forelimbs of man, bat and porpoise are homologous. & BI07 & \textbf{H} \\
BI.I10 & The wings of birds and insects are analogous and not homologous. & The wings of birds and insects are analogous and not homologous. & BI08, BI07 & \textbf{H} \\
BI.I11 & Rudimentary organs exist to complete the plan. & Interpretation: useless organs are present to complete a design. In T2 they are inherited remnants of organs useful to ancestors. & BI11$^{\circ}$, BI05a$^{\ast}$ & \textbf{U} \\
BI.I12 & Rudimentary organs are useless or nearly so. & Vestigial organs are useless or nearly so. & BI11$^{\circ}$ & \textbf{H} \\
BI.I13 & Embryos of different species in a class resemble each other more than the adults do. & Embryos of different species in a class resemble each other more than the adults do. & BI12 & \textbf{H} \\
BI.I14 & Every organism is perfectly adapted to its conditions of life. & Every organism is perfectly adapted to its conditions. [Denied: adaptation is relative to competitors and often imperfect.] & BI09$^{\circ}$ & \textbf{X} \\
BI.I15 & Organs such as the eye are fitted to their functions. & Organs such as the eye are adapted to their functions. & BI09$^{\circ}$ & \textbf{H} \\
BI.I16 & Each species arose in a single region and spread from it. & Each species arose in a single region and spread from it. & BI10b, BI01$^{\circ}$ & \textbf{H} \\
BI.I17 & Each species was created by a separate act in its centre of creation. & Interpretation: each species was brought into being by a separate act. T2 recognizes no such acts; species arise by descent. & BI10a$^{\ast}$, BI01$^{\circ}$ & \textbf{U} \\
BI.I18 & Whole faunas were destroyed by sudden catastrophes and replaced. & Whole faunas were destroyed by sudden catastrophes and replaced. [Denied in Darwin's gradualism.] & BI13$^{\circ}$ & \textbf{X} \\
BI.I19 & Fossils appear in a determinate order in successive strata. & Fossils appear in a determinate order in successive strata. & BI15 & \textbf{H} \\
BI.I20 & Extinct and living forms in the same continent are allied. & Extinct and living forms in the same continent are allied. & BI14, BI01$^{\circ}$ & \textbf{H} \\
BI.I21 & Living beings form a single linear scale without gaps. & Interpretation: organisms form one gapless linear series. T2 has a branching tree with extinct lineages. & BI16a$^{\ast}$ & \textbf{U} \\
BI.I22 & Domestic races vary more than wild species. & Domestic breeds vary more than wild species. & BI17, BI23$^{\circ}$, BI01$^{\circ}$ & \textbf{H} \\
BI.I23 & Hybrids between distinct species are generally sterile. & Hybrids between distinct species are generally sterile. & BI18$^{\circ}$, BI01$^{\circ}$ & \textbf{H} \\
BI.I24 & Hybrid sterility is a special provision to prevent the confusion of species. & Hybrid sterility is a special provision against the confusion of species. [Denied: it is incidental on other differences (Origin, ch.~8).] & BI18$^{\circ}$ & \textbf{X} \\
BI.I25 & More individuals are born than can survive. & More individuals are born than can survive. & BI19 & \textbf{H} \\
BI.I26 & Each species keeps its station in a designed balance of nature. & Each species keeps its station in a designed balance. [Denied: the relations shift as species change and go extinct.] & BI20$^{\circ}$, BI01$^{\circ}$ & \textbf{X} \\
BI.I27 & Instincts are fixed endowments, perfect for their purpose. & Instincts are fixed endowments, perfect for their purpose. [Denied: they vary and are shaped by selection.] & BI21$^{\circ}$ & \textbf{X} \\
BI.I28 & Sports occasionally appear in domestic breeds. & Large discontinuous variations occasionally appear in domestic breeds. & BI22, BI17 & \textbf{H} \\
BI.I29 & Higher organisms have more differentiated parts. & Organisms with more differentiated parts are those called higher. & BI24 & \textbf{H} \\
BI.I30 & Extinct species do not reappear. & Extinct species do not reappear. & BI13$^{\circ}$, BI01$^{\circ}$ & \textbf{H} \\
BI.I31 & The species of a genus resemble each other more than species of different genera. & The species of a genus resemble each other more than species of different genera. & BI03, BI01$^{\circ}$ & \textbf{H} \\
\bottomrule
\end{inftable}

\Needspace{10\baselineskip}\subsection*{PHYS: inferences}
\addcontentsline{toc}{subsection}{PHYS: inferences}
\begin{inftable}
GA.I01 & Blood is made continuously in the liver from the chyle of digested food. & Blood is made continuously in the liver from chyle. [Denied: the same blood circulates; Harvey's quantitative argument.] & GA06$^{\circ}$, GA16$^{\circ}$, GA01$^{\circ}$ & \textbf{X} \\
GA.I02 & Venous blood is distributed from the liver through the veins and consumed by the parts it nourishes. & Venous blood flows from the liver through the veins to the parts and is consumed there. [Denied: it flows back to the heart.] & GA01$^{\circ}$, GA24$^{\circ}$, GA19$^{\circ}$ & \textbf{X} \\
GA.I03 & Blood passes from the right to the left ventricle through pores in the septum. & Interpretation: blood crosses the septum through hidden channels. T2 recognizes none; blood passes from right to left through the lungs. & GA08a$^{\ast}$, GA07$^{\circ}$ & \textbf{U} \\
GA.I04 & The arterial vein carries blood from the right ventricle to the lungs. & The pulmonary artery carries blood from the right ventricle to the lungs. & GA09$^{\circ}$, GA20$^{\circ}$ & \textbf{H} \\
GA.I05 & The venous artery carries air from the lungs to the left ventricle. & The pulmonary vein carries air from the lungs to the left ventricle. [Denied: it carries blood.] & GA10$^{\circ}$, GA20$^{\circ}$ & \textbf{X} \\
GA.I06 & Vital spirit is elaborated in the left ventricle from air and blood. & Interpretation: an active principle is made in the left ventricle from air and blood. T2 recognizes no such principle. & GA03a$^{\ast}$, GA07$^{\circ}$ & \textbf{U} \\
GA.I07 & The arteries carry blood mixed with vital spirit to the whole body. & Interpretation: the arteries distribute blood with that principle. T2 accepts the distribution of arterial blood but not the principle. & GA02$^{\circ}$, GA03a$^{\ast}$ & \textbf{U} \\
GA.I08 & Arterial blood is brighter than venous blood. & Arterial blood is brighter than venous blood. & GA02$^{\circ}$, GA01$^{\circ}$ & \textbf{H} \\
GA.I09 & The active motion of the heart is diastole, by which it attracts blood. & Interpretation: the heart draws blood in by an active power of dilation. In T2 the active phase is systole. & GA11$^{\circ}$, GA18a$^{\ast}$ & \textbf{U} \\
GA.I10 & In systole the heart expels blood into the aorta. & In systole the heart expels blood into the aorta. & GA12, GA25 & \textbf{H} \\
GA.I11 & The arteries pulsate by a faculty conveyed from the heart through their coats. & Interpretation: the arteries pulsate by a power of their own conveyed through their walls. In T2 they are distended by the blood the heart ejects. & GA13$^{\circ}$, GA18a$^{\ast}$ & \textbf{U} \\
GA.I12 & The valves of the veins retard the flow of blood to the extremities. & The venous valves retard the flow of blood to the extremities. [Denied: they prevent flow away from the heart.] & GA17$^{\circ}$ & \textbf{X} \\
GA.I13 & The heart is the seat of the innate heat. & Interpretation: the heart is the seat of a vital heat. T2 recognizes body heat but not a vital principle seated in the heart. & GA07$^{\circ}$, GA14a$^{\ast}$ & \textbf{U} \\
GA.I14 & Respiration cools the heart and supplies the matter of vital spirit. & Interpretation: breathing cools the vital heat and supplies the matter of a vital principle. Harvey left the function of respiration open; T2 as continued explains it otherwise. & GA21$^{\circ}$, GA20$^{\circ}$, GA03a$^{\ast}$ & \textbf{U} \\
GA.I15 & Sooty vapours from the left ventricle are expelled through the venous artery to the lungs. & Interpretation: waste passes back from the left ventricle to the lungs. T2 excludes such backflow: the mitral valve prevents it. & GA15a$^{\ast}$, GA10$^{\circ}$ & \textbf{U} \\
GA.I16 & Expired air carries away waste products. & Expired air carries away waste products. & GA15b, GA21$^{\circ}$ & \textbf{H} \\
GA.I17 & Animal spirit is elaborated from vital spirit in the rete mirabile. & Interpretation: a further principle is elaborated in a network at the base of the human brain. T2 recognizes neither the principle nor the human network. & GA05a$^{\ast}$, GA26a$^{\ast}$, GA03a$^{\ast}$ & \textbf{U} \\
GA.I18 & What passes along the nerves from the brain produces motion of the muscles. & Nerve conduction from the brain produces motion of the muscles. & GA05b & \textbf{H} \\
GA.I19 & The ox has a network of arteries at the base of its brain. & The ox has a network of arteries at the base of its brain (the carotid rete). & GA26b & \textbf{H} \\
GA.I20 & Natural spirit is elaborated in the liver and carried by the veins. & Interpretation: a principle of nutrition is elaborated in the liver. T2 recognizes no such principle. & GA04a$^{\ast}$, GA06$^{\circ}$ & \textbf{U} \\
GA.I21 & Arteries and veins communicate at their extremities. & Arteries and veins communicate at their extremities (through capillaries, seen by Malpighi in 1661). & GA22 & \textbf{H} \\
GA.I22 & The portal vein carries chyle from the intestines to the liver. & The hepatic portal vein carries chyle to the liver. [Denied: it carries blood with absorbed nutrients; the chyle passes by the lacteals.] & GA23$^{\circ}$, GA16$^{\circ}$, GA06$^{\circ}$ & \textbf{X} \\
GA.I23 & The aorta carries arterial blood from the left ventricle to the body. & The aorta carries arterial blood from the left ventricle to the body. & GA25, GA02$^{\circ}$ & \textbf{H} \\
GA.I24 & A ligature on an artery stops the pulse beyond it. & A ligature on an artery stops the pulse beyond it. & GA13$^{\circ}$ & \textbf{H} \\
GA.I25 & Each part attracts the nourishment it needs by its attractive faculty. & Interpretation: each part draws nourishment by a power of its own. T2 recognizes no such power. & GA18a$^{\ast}$, GA19$^{\circ}$ & \textbf{U} \\
GA.I26 & The parts are nourished by matter carried to them in the blood. & The parts are nourished by matter carried to them in the blood. & GA19$^{\circ}$, GA01$^{\circ}$ & \textbf{H} \\
GA.I27 & The lungs are nourished by blood brought by the arterial vein. & The lungs are nourished by blood brought by the pulmonary artery. [Denied: that blood passes through the lungs to the left heart.] & GA09$^{\circ}$, GA20$^{\circ}$, GA19$^{\circ}$ & \textbf{X} \\
GA.I28 & The living body is warm, and grows cold at death. & The living body is warm, and grows cold at death. & GA14b & \textbf{H} \\
GA.I29 & Blood flows out of a cut vein and a cut artery. & Blood flows out of a cut vein and a cut artery. & GA01$^{\circ}$, GA02$^{\circ}$ & \textbf{H} \\
GA.I30 & In systole the heart becomes harder and smaller. & In systole the heart becomes harder and smaller. & GA12, GA07$^{\circ}$ & \textbf{H} \\
\bottomrule
\end{inftable}

\Needspace{10\baselineskip}\subsection*{OPTI: inferences}
\addcontentsline{toc}{subsection}{OPTI: inferences}
\begin{inftable}
OP.I01 & Light consists of corpuscles emitted by luminous bodies. & Interpretation: light is a stream of small bodies. The wave theory recognizes none. & OP02a$^{\ast}$ & \textbf{U} \\
OP.I02 & Rays travel in straight lines in a uniform medium. & Light travels in straight lines in a uniform medium. [Holds when apertures and obstacles are large compared with the wavelength.] & OP01$^{\circ}$ & \textbf{R} \\
OP.I03 & Refraction toward the normal arises from the attraction of the corpuscles by the denser medium. & Interpretation: refraction is due to an attraction of light-particles by the denser medium. In T2 it is due to the lower wave speed in that medium. & OP02a$^{\ast}$, OP05$^{\circ}$ & \textbf{U} \\
OP.I04 & Light travels faster in water than in air. & Light travels faster in water than in air. [Denied: slower, as Foucault found in 1850.] & OP06$^{\circ}$ & \textbf{X} \\
OP.I05 & Sunlight is a mixture of rays of different refrangibility. & Sunlight is a superposition of light of different wavelengths, refracted differently. & OP13, OP05$^{\circ}$, OP01$^{\circ}$ & \textbf{H} \\
OP.I06 & Homogeneal light keeps its colour and refrangibility under reflection and refraction. & Monochromatic light keeps its wavelength in reflection and refraction (its frequency is unchanged). & OP14, OP04, OP05$^{\circ}$ & \textbf{H} \\
OP.I07 & Each corpuscle passes alternately into fits of easy reflection and of easy transmission. & Interpretation: each light-particle alternates between two dispositions. In T2 the periodicity belongs to the wave, not to particles. & OP03a$^{\ast}$, OP02a$^{\ast}$ & \textbf{U} \\
OP.I08 & The thicknesses of the film at the bright rings are as the odd numbers 1, 3, 5, \ldots & The film thicknesses at the bright rings in reflected light are as the odd numbers 1, 3, 5 (in quarter-wavelengths). & OP11, OP03b & \textbf{H} \\
OP.I09 & The interval of the fits is greater for red than for violet light. & The wavelength is greater for red than for violet light. & OP03b, OP04 & \textbf{H} \\
OP.I10 & Light passing near an edge is inflected and forms coloured fringes outside the shadow. & Light diffracted near an edge forms coloured fringes outside the shadow. & OP07$^{\circ}$, OP16$^{\circ}$ & \textbf{H} \\
OP.I11 & Inflection is caused by forces of the edge acting on the corpuscles. & Interpretation: the edge exerts forces on passing light-particles. In T2 fringes arise from the interference of waves. & OP07$^{\circ}$, OP02a$^{\ast}$ & \textbf{U} \\
OP.I12 & Light does not bend into the geometrical shadow. & Light does not enter the geometrical shadow. [Denied: the bright spot at the centre of a disk's shadow.] & OP16$^{\circ}$, OP01$^{\circ}$ & \textbf{X} \\
OP.I13 & The sides of a ray determine its refraction in a second Iceland crystal. & The polarization of a beam determines its refraction in a second calcite crystal. & OP08b, OP10 & \textbf{H} \\
OP.I14 & The sides are original properties of the corpuscles. & Interpretation: the sides are properties of the particles. In T2 polarization is the orientation of a transverse vibration. & OP08a$^{\ast}$ & \textbf{U} \\
OP.I15 & Light reflected from glass at a certain angle is polarized. & Light reflected from glass at a certain angle is polarized (Brewster's angle). & OP09$^{\circ}$, OP15$^{\circ}$ & \textbf{H} \\
OP.I16 & Iceland crystal divides a beam into two beams polarized at right angles. & Calcite divides a beam into two beams polarized at right angles. & OP10, OP09$^{\circ}$ & \textbf{H} \\
OP.I17 & Polarization consists in the alignment of the axes of the light particles. & Interpretation: polarization aligns the axes of light-particles. T2 has no particles to align. & OP09$^{\circ}$, OP02a$^{\ast}$ & \textbf{U} \\
OP.I18 & Starlight is aberrated by the ratio of the Earth's speed to the speed of light. & Starlight is aberrated by the ratio of the Earth's speed to the speed of light. & OP12 & \textbf{H} \\
OP.I19 & Intensity diminishes as the inverse square of the distance. & Intensity diminishes as the inverse square of the distance. & OP17 & \textbf{H} \\
OP.I20 & Reflection is caused by a repulsive power acting at a distance from the surface. & Interpretation: a repulsive force near the surface turns the particles back. In T2 reflection follows from the boundary conditions on the wave. & OP15$^{\circ}$, OP02a$^{\ast}$ & \textbf{U} \\
OP.I21 & The angle of reflection equals the angle of incidence. & The angle of reflection equals the angle of incidence. & OP15$^{\circ}$, OP01$^{\circ}$ & \textbf{H} \\
OP.I22 & For homogeneal light the sine of incidence bears a constant ratio to the sine of refraction. & For monochromatic light the ratio of the sines of incidence and refraction is constant. & OP05$^{\circ}$, OP14, OP01$^{\circ}$ & \textbf{H} \\
OP.I23 & A prism spreads sunlight into a spectrum from red, least refracted, to violet, most refracted. & A prism disperses sunlight into a spectrum from long wavelengths, least refracted, to short, most refracted. & OP18, OP13, OP04 & \textbf{H} \\
OP.I24 & The rings in transmitted light are complementary in colour to those in reflected light. & The rings in transmitted light are complementary to those in reflected light. & OP11, OP04 & \textbf{H} \\
OP.I25 & The colours of thin plates depend on the thickness of the plate. & The colours of thin films depend on their thickness through the path difference of the reflected waves. & OP11, OP03b & \textbf{H} \\
\bottomrule
\end{inftable}

\Needspace{10\baselineskip}\subsection*{MIAS: inferences}
\addcontentsline{toc}{subsection}{MIAS: inferences}
\begin{inftable}
MI.I01 & Epidemic diseases are caused by miasmata from decaying organic matter. & Interpretation: a poisonous emanation of decay, breathed in, causes epidemics. T2 recognizes no such single poison; epidemics are caused by specific organisms. & MI01a$^{\ast}$ & \textbf{U} \\
MI.I02 & Intermittent fevers are most prevalent in marshy districts and diminish when the marshes are drained. & Malaria is most prevalent in marshy districts and diminishes when the marshes are drained (which removes the breeding places of the mosquitoes that carry Plasmodium). & MI09, MI01b & \textbf{H} \\
MI.I03 & Zymotic diseases are produced by a process in the blood analogous to fermentation. & Interpretation: a chemical fermentation in the blood produces these diseases. In T2 they are produced by the multiplication of organisms in the body. & MI02$^{\circ}$, MI03c$^{\ast}$ & \textbf{U} \\
MI.I04 & Smallpox is communicated by contact with the sick or their clothing. & Smallpox is transmitted by contact with the sick or their clothing. & MI04, MI05a & \textbf{H} \\
MI.I05 & Contagia are living organisms. & Pathogenic agents are living organisms. & MI05b & \textbf{H} \\
MI.I06 & Contagia are chemical poisons that multiply by inducing decomposition. & Interpretation: the agents are chemical poisons multiplying by catalysis. T2 has no such poisons; the agents reproduce as organisms. & MI05c$^{\ast}$ & \textbf{U} \\
MI.I07 & Yeast is necessary for alcoholic fermentation. & Yeast cells are necessary for alcoholic fermentation. & MI03a, MI20$^{\circ}$ & \textbf{H} \\
MI.I08 & Fermentation is a decomposition propagated from a decomposing body. & Interpretation: fermentation is decomposition propagated from a decomposing body. In T2 it is the metabolism of living cells or the action of their enzymes. & MI20$^{\circ}$, MI03c$^{\ast}$ & \textbf{U} \\
MI.I09 & Putrefaction arises in dead organic matter from the action of air alone. & Microbial decomposition arises in dead organic matter from the action of air alone. [Denied: it requires organisms, which air may carry; sterile matter in filtered air does not putrefy.] & MI06$^{\circ}$ & \textbf{X} \\
MI.I10 & Microscopic organisms arise in infusions without parents. & Interpretation: organisms arise without parents. T2 (Pasteur, 1862) denies that they do. & MI07a$^{\ast}$ & \textbf{U} \\
MI.I11 & Boiled hay infusions can teem with organisms after cooling. & Boiled hay infusions can teem with organisms after cooling (grown from spores that survive boiling). & MI07b & \textbf{H} \\
MI.I12 & Epidemics follow constitutions of the atmosphere. & Interpretation: epidemics follow states of the air. T2 recognizes no such states as causes, though weather affects vectors and transmission. & MI08a$^{\ast}$ & \textbf{U} \\
MI.I13 & Removal of filth reduces the mortality of epidemic disease. & Removal of organic refuse reduces the mortality of infectious disease. & MI11$^{\circ}$ & \textbf{H} \\
MI.I14 & Disease arises from breathing the emanations of filth. & Interpretation: the smell of filth causes disease. T2 denies that smells cause disease; filth harbours pathogens. & MI11$^{\circ}$, MI01a$^{\ast}$ & \textbf{U} \\
MI.I15 & Sewer gas causes typhoid. & Interpretation: a gas from drains causes typhoid. In T2 typhoid is caused by Salmonella Typhi, transmitted by contaminated water and food. & MI12a$^{\ast}$ & \textbf{U} \\
MI.I16 & Cholera is communicated by morbid matter in the evacuations, swallowed in water. & Cholera is transmitted by Vibrio cholerae in the evacuations, swallowed in water. & MI13a, MI21$^{\circ}$ & \textbf{H} \\
MI.I17 & The cholera poison must mature in porous soil before it causes disease. & Interpretation: the cholera agent must ripen in the soil. T2 recognizes no such ripening. & MI13b$^{\ast}$ & \textbf{U} \\
MI.I18 & Cholera mortality in London in 1849 was lower in districts of higher elevation. & Cholera mortality in London in 1849 was lower in districts of higher elevation. & MI21$^{\circ}$ & \textbf{H} \\
MI.I31 & The cholera miasma was most concentrated in the low-lying districts by the Thames. & Vibrio cholerae was most prevalent in the low-lying districts by the Thames. [Holds for London in 1849, where those districts drew contaminated water.] & MI01c, MI21$^{\circ}$ & \textbf{R} \\
MI.I19 & Agents that destroy offensive smells prevent disease by destroying miasma. & Interpretation: deodorants prevent disease by destroying an emanation. In T2 some deodorants are disinfectants and prevent disease by killing microorganisms. & MI14$^{\circ}$, MI01a$^{\ast}$ & \textbf{U} \\
MI.I20 & Washing the hands in chloride of lime reduced mortality from puerperal fever. & Washing the hands in hypochlorite reduced mortality from puerperal sepsis. & MI14$^{\circ}$, MI17$^{\circ}$ & \textbf{H} \\
MI.I21 & Puerperal fever arises from an epidemic influence in the ward atmosphere. & Interpretation: an epidemic influence of the air causes puerperal fever. In T2 it is carried by the hands and instruments of attendants. & MI17$^{\circ}$, MI08a$^{\ast}$ & \textbf{U} \\
MI.I22 & Vaccination with cowpox protects against smallpox. & Vaccination with cowpox protects against smallpox. & MI15 & \textbf{H} \\
MI.I23 & Hospital gangrene is more common in crowded wards. & Necrotizing wound infection is more common in crowded wards. & MI16 & \textbf{H} \\
MI.I24 & Each zymotic disease breeds true: smallpox produces smallpox and never another disease. & Each infectious disease breeds true: variola produces smallpox and never another disease. & MI02$^{\circ}$, MI05a & \textbf{H} \\
MI.I25 & Continued fever is a single disease. & Typhus and typhoid are one disease. [Denied: distinguished clinically by William Jenner (1849--1851), aetiologically later.] & MI10$^{\circ}$ & \textbf{X} \\
MI.I26 & Poverty and intemperance predispose to epidemic disease. & Poverty and intemperance increase susceptibility to infectious disease. & MI19 & \textbf{H} \\
MI.I27 & Infection can be carried by clothes and bedding. & Infection can be carried by clothes and bedding. & MI18, MI04 & \textbf{H} \\
MI.I28 & Intermittent fever is cured by cinchona bark. & Malaria is cured by cinchona bark (quinine). & MI09 & \textbf{H} \\
MI.I29 & Diastase converts starch into sugar. & Amylase converts starch into sugar. & MI03b & \textbf{H} \\
MI.I30 & Putrid matter introduced into the blood produces fever. & Putrid matter introduced into the blood produces fever (by infection). & MI06$^{\circ}$, MI18 & \textbf{H} \\
\bottomrule
\end{inftable}

\Needspace{10\baselineskip}\subsection*{HERE: inferences}
\addcontentsline{toc}{subsection}{HERE: inferences}
\begin{inftable}
HR.I01 & Every cell throws off gemmules, which collect in the reproductive elements. & Interpretation: particles shed by every cell collect in the gametes. T2 has no such particles; the hereditary material is in the chromosomes of the germ line. & HR01a$^{\ast}$, HR17$^{\circ}$ & \textbf{U} \\
HR.I02 & The offspring blends the characters of the parents, and the blend is not undone in later generations. & Quantitative characters of offspring are intermediate, and the intermediacy persists in later generations. [Denied: segregation restores parental types.] & HR02$^{\circ}$, HR11$^{\circ}$ & \textbf{X} \\
HR.I03 & A single favourable variation is swamped by crossing within a few generations. & A rare favourable variant is diluted to nothing by crossing within a few generations. [Denied: Mendelian alleles are not diluted.] & HR02$^{\circ}$, HR10$^{\circ}$ & \textbf{X} \\
HR.I04 & In crosses one parent is often prepotent in transmitting a character. & In crosses one allele is often dominant in its effect on a trait. & HR03, HR11$^{\circ}$ & \textbf{H} \\
HR.I05 & Crossing distinct pure breeds of pigeon can produce the colouring of the rock pigeon. & Crossing distinct homozygous breeds of pigeon can produce the rock-pigeon phenotype (by recombination). & HR04$^{\circ}$, HR07$^{\circ}$ & \textbf{H} \\
HR.I06 & Reversion is due to dormant gemmules transmitted through generations. & Reversion is due to recessive alleles transmitted unexpressed through generations. [Holds for reversion by recessives; reversions by recombination are explained otherwise.] & HR04$^{\circ}$, HR01b & \textbf{R} \\
HR.I07 & Characters acquired by use or disuse are inherited. & Interpretation: effects of use and disuse are transmitted. T2 (with Weismann) recognizes no mechanism for this. & HR05a$^{\ast}$, HR11$^{\circ}$ & \textbf{U} \\
HR.I08 & A previous sire influences the offspring of later matings. & Interpretation: a previous sire affects later offspring. T2 recognizes no route for such influence. & HR06a$^{\ast}$ & \textbf{U} \\
HR.I09 & Pure breeds breed true. & Homozygous lines breed true. & HR07$^{\circ}$ & \textbf{H} \\
HR.I10 & The offspring of hybrids are variable, some resembling each grandparent. & The offspring of heterozygotes are variable, some resembling each grandparent (segregation). & HR08$^{\circ}$, HR11$^{\circ}$ & \textbf{H} \\
HR.I11 & The offspring of hybrids remain uniformly intermediate. & The offspring of heterozygotes are uniformly intermediate. [Denied: they segregate.] & HR08$^{\circ}$, HR02$^{\circ}$ & \textbf{X} \\
HR.I12 & Sports are sometimes inherited. & Mutations are sometimes inherited. & HR09 & \textbf{H} \\
HR.I13 & Individual variation is heritable in some degree. & Individual variation is heritable in some degree. [Holds on the reading heritable variation; fails for variation induced by the environment.] & HR10$^{\circ}$ & \textbf{R} \\
HR.I14 & Selection within a pure breed shifts its mean. & Selection within a homozygous line shifts its mean. [Denied: Johannsen's pure lines (1903).] & HR10$^{\circ}$, HR07$^{\circ}$ & \textbf{X} \\
HR.I15 & Each offspring receives hereditary tendencies from both parents. & Each offspring receives its genotype from both parents. & HR12, HR17$^{\circ}$ & \textbf{H} \\
HR.I16 & The offspring of exceptional parents regress toward the mean of the race. & The offspring of exceptional parents regress toward the population mean. & HR14 & \textbf{H} \\
HR.I17 & Parents contribute one half of the heritage, grandparents one quarter, and so on. & Interpretation: the heritage is composed of fixed fractions from each ancestor. T2 has no such fractions; each parent transmits one allele at each locus. & HR13a$^{\ast}$ & \textbf{U} \\
HR.I18 & Filial stature regresses on mid-parental stature with a coefficient of about two thirds. & Filial stature regresses on mid-parental stature with a coefficient of about two thirds. & HR13b, HR14 & \textbf{H} \\
HR.I19 & The hereditary substance is continuous and is not affected by the body. & The germ-line hereditary material is continuous and is not altered by the body. & HR15 & \textbf{H} \\
HR.I20 & A character can lie latent for a generation and reappear. & A recessive trait can skip a generation and reappear. & HR16, HR11$^{\circ}$ & \textbf{H} \\
HR.I21 & Each parent contributes one reproductive element to the offspring. & Each parent contributes one gamete to the offspring. & HR17$^{\circ}$ & \textbf{H} \\
HR.I22 & The intermediate stature of children of tall and short parents is due to blending. & The intermediate stature of children of tall and short parents is due to quantitative inheritance. [Holds: many factors with additive effect give intermediate means.] & HR02$^{\circ}$, HR11$^{\circ}$ & \textbf{R} \\
HR.I23 & A dilute rare variation in a large population cannot be raised in frequency by selection. & A rare variant in a large population cannot be raised in frequency by selection. [Denied: selection raises the frequency of an undiluted allele.] & HR02$^{\circ}$, HR10$^{\circ}$ & \textbf{X} \\
\bottomrule
\end{inftable}

